\documentclass[11pt,a4paper]{article}
% preamble.tex — estilo común de todos los apuntes Froth.
% Cada apunte hace % preamble.tex — estilo común de todos los apuntes Froth.
% Cada apunte hace % preamble.tex — estilo común de todos los apuntes Froth.
% Cada apunte hace \input{preamble} justo después de \documentclass.
\usepackage[margin=2.4cm]{geometry}
\usepackage{xcolor}
\definecolor{litblue}{RGB}{31,79,121}
\definecolor{litgreen}{RGB}{27,94,32}
\definecolor{litorange}{RGB}{230,81,0}
\usepackage{parskip}
\usepackage{titlesec}
\titleformat{\section}{\Large\bfseries\color{litblue}}{\thesection}{0.6em}{}
\titleformat{\subsection}{\large\bfseries\color{litblue}}{\thesubsection}{0.6em}{}
\usepackage{tcolorbox}
\usepackage{fancyhdr}
\pagestyle{fancy}
\fancyhf{}
\fancyhead[L]{\small\color{litblue}Froth — Apuntes de ML}
\fancyhead[R]{\small\thepage}
\renewcommand{\headrulewidth}{0.4pt}
\usepackage[colorlinks=true,linkcolor=litblue,urlcolor=litblue]{hyperref}

% Cabecera de cada apunte: \cabecera{numero}{titulo}
\newcommand{\cabecera}[2]{%
  \begin{tcolorbox}[colback=litblue,coltext=white,colframe=litblue,arc=2mm]
    {\Large\bfseries Apunte #1 — #2}\\[3pt]
    {\small Proyecto Froth · aprender Machine Learning construyendo}
  \end{tcolorbox}\vspace{0.8em}}

% Cajas temáticas
\newtcolorbox{concepto}[1]{colback=blue!4,colframe=litblue,title={Concepto: #1},arc=1.5mm}
\newtcolorbox{practica}{colback=green!5,colframe=litgreen,title={Pruébalo tú (teclado en mano)},arc=1.5mm}
\newtcolorbox{ojo}{colback=orange!8,colframe=litorange,title={Ojo / error típico},arc=1.5mm}
\newtcolorbox{resumen}{colback=gray!8,colframe=gray!60,title={En una frase},arc=1.5mm}

\newcommand{\codigo}[1]{\texttt{#1}}
 justo después de \documentclass.
\usepackage[margin=2.4cm]{geometry}
\usepackage{xcolor}
\definecolor{litblue}{RGB}{31,79,121}
\definecolor{litgreen}{RGB}{27,94,32}
\definecolor{litorange}{RGB}{230,81,0}
\usepackage{parskip}
\usepackage{titlesec}
\titleformat{\section}{\Large\bfseries\color{litblue}}{\thesection}{0.6em}{}
\titleformat{\subsection}{\large\bfseries\color{litblue}}{\thesubsection}{0.6em}{}
\usepackage{tcolorbox}
\usepackage{fancyhdr}
\pagestyle{fancy}
\fancyhf{}
\fancyhead[L]{\small\color{litblue}Froth — Apuntes de ML}
\fancyhead[R]{\small\thepage}
\renewcommand{\headrulewidth}{0.4pt}
\usepackage[colorlinks=true,linkcolor=litblue,urlcolor=litblue]{hyperref}

% Cabecera de cada apunte: \cabecera{numero}{titulo}
\newcommand{\cabecera}[2]{%
  \begin{tcolorbox}[colback=litblue,coltext=white,colframe=litblue,arc=2mm]
    {\Large\bfseries Apunte #1 — #2}\\[3pt]
    {\small Proyecto Froth · aprender Machine Learning construyendo}
  \end{tcolorbox}\vspace{0.8em}}

% Cajas temáticas
\newtcolorbox{concepto}[1]{colback=blue!4,colframe=litblue,title={Concepto: #1},arc=1.5mm}
\newtcolorbox{practica}{colback=green!5,colframe=litgreen,title={Pruébalo tú (teclado en mano)},arc=1.5mm}
\newtcolorbox{ojo}{colback=orange!8,colframe=litorange,title={Ojo / error típico},arc=1.5mm}
\newtcolorbox{resumen}{colback=gray!8,colframe=gray!60,title={En una frase},arc=1.5mm}

\newcommand{\codigo}[1]{\texttt{#1}}
 justo después de \documentclass.
\usepackage[margin=2.4cm]{geometry}
\usepackage{xcolor}
\definecolor{litblue}{RGB}{31,79,121}
\definecolor{litgreen}{RGB}{27,94,32}
\definecolor{litorange}{RGB}{230,81,0}
\usepackage{parskip}
\usepackage{titlesec}
\titleformat{\section}{\Large\bfseries\color{litblue}}{\thesection}{0.6em}{}
\titleformat{\subsection}{\large\bfseries\color{litblue}}{\thesubsection}{0.6em}{}
\usepackage{tcolorbox}
\usepackage{fancyhdr}
\pagestyle{fancy}
\fancyhf{}
\fancyhead[L]{\small\color{litblue}Froth — Apuntes de ML}
\fancyhead[R]{\small\thepage}
\renewcommand{\headrulewidth}{0.4pt}
\usepackage[colorlinks=true,linkcolor=litblue,urlcolor=litblue]{hyperref}

% Cabecera de cada apunte: \cabecera{numero}{titulo}
\newcommand{\cabecera}[2]{%
  \begin{tcolorbox}[colback=litblue,coltext=white,colframe=litblue,arc=2mm]
    {\Large\bfseries Apunte #1 — #2}\\[3pt]
    {\small Proyecto Froth · aprender Machine Learning construyendo}
  \end{tcolorbox}\vspace{0.8em}}

% Cajas temáticas
\newtcolorbox{concepto}[1]{colback=blue!4,colframe=litblue,title={Concepto: #1},arc=1.5mm}
\newtcolorbox{practica}{colback=green!5,colframe=litgreen,title={Pruébalo tú (teclado en mano)},arc=1.5mm}
\newtcolorbox{ojo}{colback=orange!8,colframe=litorange,title={Ojo / error típico},arc=1.5mm}
\newtcolorbox{resumen}{colback=gray!8,colframe=gray!60,title={En una frase},arc=1.5mm}

\newcommand{\codigo}[1]{\texttt{#1}}

\begin{document}
\cabecera{41}{El veredicto del re-fine-tune: más datos NO mejoró}

\section{La hipótesis (tu duda del learning log)}

El fine-tune v1 (specter-mineral, 1.290 papers) ganaba 14/22 en DBCV pero castigaba la
periferia (hasta $-0.43$). Hipótesis natural: con \textbf{3$\times$ más datos}
(re-cosecha con keys $\to$ 4.045 papers) el modelo v2 debería generalizar mejor y dejar
de penalizar los topics lejanos. Lo pusimos a prueba, medido.

\section{El experimento limpio}

Entrenamos v2 (4.045 pares título$\leftrightarrow$abstract, MNRL, 249s en tu RTX 5060) y
comparamos \textbf{los tres modelos sobre los MISMOS 22 corpus nuevos} por DBCV: base
(SPECTER genérico), v1 (fine-tune viejo) y v2 (fine-tune grande).

\begin{center}
\begin{tabular}{lccc}
 & gana a base & delta medio \\
\hline
v1 (1.290 papers) & \textbf{16/22} & \textbf{+0.093} \\
v2 (4.045 papers) & 14/22 & +0.069 \\
v2 vs v1 directo & 10/22 & \textbf{$-$0.024} \\
\end{tabular}
\end{center}

\begin{concepto}{el resultado: escalar datos NO ayudó (y hasta empeoró un poco)}
v2 gana a base MENOS veces que v1 (14 vs 16) y con menor margen. Cabeza a cabeza,
v2 le gana a v1 en solo 10/22 con delta medio negativo. \textbf{El fine-tune pequeño y
curado fue igual o mejor que el grande.}
\end{concepto}

\section{Por qué (la lección de verdad)}

La re-cosecha 3$\times$ trajo más papers, pero también más DIVERSOS y RUIDOSOS
(colisiones de términos multi-dominio como ``attrition''/``autogenous'', más subdominios).
El fine-tuning auto-supervisado arrastra el modelo hacia el \emph{centro de masa} del
corpus de entrenamiento (apunte 31); al agrandar el corpus con ruido, ese centro se
mueve hacia lo genérico y la especialización se DILUYE.

\begin{ojo}
En fine-tuning auto-supervisado, la CALIDAD y coherencia del corpus importa más que la
cantidad. ``Más datos'' es un reflejo, no una ley. La varianza sigue altísima: v2 va de
$+0.47$ (escorias XCT) a $-0.51$ (finos de silicato de hierro) --- el fine-tune es una
apuesta por-topic, no una mejora uniforme.
\end{ojo}

\section{Qué queda en producción}

Nada cambia: producción sigue con el SPECTER \textbf{base} (el más estable across topics).
Los fine-tunes son el RESULTADO DE INVESTIGACIÓN. v1 se conserva
(\texttt{models/specter-mineral-v1}); v2 queda como referencia. Evidencia completa en
\texttt{2\_Datos/model\_comparison\_3way.csv}.

\begin{resumen}
Re-fine-tune con 3$\times$ datos: v2 gana a base 14/22 (v1 ganaba 16/22), y pierde contra
v1 cabeza a cabeza (delta medio $-0.024$). Más datos NO mejoró --- el corpus grande y
ruidoso diluyó la especialización. Lección de tesis honesta: en fine-tuning
auto-supervisado, la curación del corpus $>$ el tamaño. ``No mejora'' también es
resultado.
\end{resumen}

\end{document}
